\chapter{\(3\) feuillets temporels, calcul nilpotent et connexion tritifiée de Cartan}
\label{chap:tres-hojas-variacion}

\section{Famille quadratique de TRIT}

Pour \(\tau\in\{+1,0,-1\}\), considérons l'algèbre réelle

\[
\mathbb A_\tau=\mathbb R[J_\tau]/(J_\tau^2+\tau),
\qquad
J_\tau^2=-\tau I.
\]

L'exponentielle se calcule exactement :

\[
e^{tJ_{+1}}=\cos t+J_{+1}\sin t,
\]

\[
e^{tJ_0}=I+tJ_0,
\]

\[
e^{tJ_{-1}}=\cosh t+J_{-1}\sinh t.
\]

Les trois feuillets possèdent des fonctions géométriques différenciées :

\begin{center}
\begin{tabular}{ccl}
\toprule
\(\tau\)&\(J_\tau^2\)&r\'egimen\\
\midrule
\(+1\)&\(-I\)&elliptique : clôture, retour et mémoire\\
\(0\)&\(0\)&parabolique : tangence, variation et présent\\
\(-1\)&\(+I\)&hyperbolique : ouverture, boost et propagation\\
\bottomrule
\end{tabular}
\end{center}

Soit \(\widehat\tau=P_+-P_-\), avec \(P_0=I-P_+-P_-\). Alors

\[
\widehat\tau^3=\widehat\tau.
\]

Une involution qui échange les feuillets extrêmes et conserve le feuillet central
satisfait

\[
\mathsf J\widehat\tau\mathsf J^{-1}=-\widehat\tau.
\]

Le seuil reste fixe et les orientations de fermeture et d'ouverture
s'échangent.

\section{Le feuillet nul comme 1.er jet}

Dans \(\mathbb A_0\), écrivons \(\epsilon=J_0\), avec \(\epsilon^2=0\). Pour
toute fonctionnelle différentiable \(F\),

\begin{equation}
\boxed{F(x+\epsilon v)=F(x)+\epsilon\,dF_x(v).}
\label{eq:rm-dual-first-jet}
\end{equation}

L'égalité est exacte car tous les termes d'ordre au moins deux
contiennent \(\epsilon^2\). En particulier, pour une action
\(S[e,\omega,\Psi]\),

\begin{equation}
S[e+\epsilon\delta e,\omega+\epsilon\delta\omega,
  \Psi+\epsilon\delta\Psi]
=S[e,\omega,\Psi]+\epsilon\,\delta S.
\label{eq:rm-action-first-jet}
\end{equation}

Le coefficient nilpotent contient les équations d'Euler--Lagrange,
l'équation de Cartan, le potentiel symplectique et les termes de bord.
Le feuillet central est donc la résidence algébrique du principe
variationnel.

\begin{theorem}[Stationnarité nilpotente]
Pour une configuration \(z\), sont équivalents :
\[
dS_z(\delta z)=0\quad\text{pour toute }\delta z,
\]
et
\[
S(z+\epsilon\delta z)=S(z)
\quad\text{dans }\mathbb A_0
\quad\text{pour toute }\delta z.
\]
\end{theorem}

\begin{proof}
L'équivalence est l'égalité des coefficients de \(1\) et \(\epsilon\) dans
\eqref{eq:rm-dual-first-jet}.
\end{proof}

Le Hamiltonien global, l'action effective du feuillet central et le
Hamiltonien modulaire du bord sont des objets coordonnés :

\[
H_{\rm glob},
\qquad
L_{0,\rm eff},
\qquad
K_{\rm HHM}=-\log\rho_{\partial}.
\]

Le premier engendre l'évolution ; le deuxième publie la variation locale ; le
troisième conserve la provenance informationnelle.

\section{Réduction de Feshbach--Schur}

Soit

\[
\mathcal H=\mathcal H_0\oplus\mathcal H_{\rm ext},
\qquad
H=
\begin{pmatrix}
H_{00}&V\\ V^*&H_{\rm ext}
\end{pmatrix}.
\]

Pour \(z\) dans l'ensemble résolvant de \(H_{\rm ext}\), la compression de la
résolvante est

\begin{equation}
P_0(z-H)^{-1}P_0
=\left[z-H_{00}-\Sigma_0(z)\right]^{-1},
\qquad
\Sigma_0(z)=V(z-H_{\rm ext})^{-1}V^*.
\label{eq:rm-feshbach-schur}
\end{equation}

L'opérateur \(\Sigma_0\) transporte vers le feuillet présent la mémoire des feuillets
elliptique et hyperbolique. Sa transformée temporelle engendre un noyau de
mémoire

\[
S_{\rm eff}[\psi_0]
=\int\psi_0^\dagger(i\hbar\partial_t-H_{00})\psi_0\,\mathrm dt
-\iint\psi_0^\dagger(t)K_0(t,t')\psi_0(t')\,\mathrm dt\,\mathrm dt'.
\]

Un développement régulier à basse fréquence,

\[
\Sigma_0(\omega)
=\Sigma_0^{(0)}+\omega\Sigma_0^{(1)}
+\omega^2\Sigma_0^{(2)}+\cdots,
\]

produit une action locale effective ; une structure spectrale singulière conserve
une mémoire non locale. Le présent reçoit ainsi la réponse de l'état global sans
introduire un temps extérieur.

\section{Connexion universelle de Cartan}

Soit \(\mathfrak g_\tau\) l'algèbre de Cartan engendrée par
\(J_{ab}=-J_{ba}\) et \(P_a^{(\tau)}\), avec les relations

\begin{align}
 [J_{ab},J_{cd}]&=
 \eta_{bc}J_{ad}-\eta_{ac}J_{bd}
 -\eta_{bd}J_{ac}+\eta_{ad}J_{bc},\\
 [J_{ab},P_c^{(\tau)}]&=
 \eta_{bc}P_a^{(\tau)}-\eta_{ac}P_b^{(\tau)},\\
 [P_a^{(\tau)},P_b^{(\tau)}]&=-\tau J_{ab}.
 \label{eq:rm-cartan-brackets}
\end{align}

Pour \(\tau=0\), on obtient l'algèbre de Poincaré ; les signes extrêmes
produisent les deux formes de courbure constante. Soient \(\omega\) une connexion
de Lorentz, \(e=e^aP_a^{(\tau)}\) une cotétrade à valeurs dans le secteur des
translations et \(\ell\) une échelle de longueur. Nous définissons la connexion
tritifiée

\begin{equation}
\mathcal A_\tau^{\rm Cart}
=\frac12\omega^{ab}J_{ab}+\ell^{-1}e^aP_a^{(\tau)}.
\label{eq:rm-trit-cartan-connection}
\end{equation}

Sa courbure est

\begin{equation}
\boxed{
\mathcal F_\tau^{\rm Cart}
=\frac12\left(F_\omega^{ab}-\tau\ell^{-2}e^a\wedge e^b\right)J_{ab}
+\ell^{-1}T^aP_a^{(\tau)},
\qquad T^a=D_\omega e^a.}
\label{eq:rm-trit-cartan-curvature}
\end{equation}

\begin{proof}
En développant
\(\mathrm d\mathcal A_\tau^{\rm Cart}
+\mathcal A_\tau^{\rm Cart}\wedge\mathcal A_\tau^{\rm Cart}\), le premier
crochet de \cref{eq:rm-cartan-brackets} produit \(F_\omega\), le deuxième
produit \(D_\omega e\) et le troisième apporte
\(-\tau\ell^{-2}e^a\wedge e^bJ_{ab}/2\).
\end{proof}

Dans le feuillet central,

\[
\mathcal F_0^{\rm Cart}
=\frac12F_\omega^{ab}J_{ab}+\ell^{-1}T^aP_a^{(0)}.
\]

La torsion est le coefficient nilpotent de la courbure. Cette égalité unit
l'interprétation variationnelle du feuillet \(0\) à son interprétation
géométrique de Cartan.

\section{Couronne nonadique et énergie torsionnelle}

Sur les arêtes de la couronne de neuf, soit

\[
(D_{9,m}f)(a)=f(t(a))-U_a f(s(a)),
\]

où \(U_a\) transporte l'orientation et la retenue. Pour un poids positif \(W_m\),

\begin{equation}
H_m=H_{{\rm clk},m}+D_{9,m}^*W_mD_{9,m},
\qquad
\mathcal E_{{\rm tor},m}(f)
=\|W_m^{1/2}D_{9,m}f\|^2.
\label{eq:rm-nonadic-torsion-energy}
\end{equation}

L'holonomie unitaire conserve la positivité. Les carrés de raffinement

\[
D_{9,m+1}J_m=K_mD_{9,m},
\qquad
K_m^*W_{m+1}K_m=W_m
\]

garantissent la compatibilité de l'énergie à toute profondeur.

La clôture ultérieure construira la soudure d'écran qui transforme ce
défaut en une réponse de Cartan. La présente section fixe le domaine
opératoire et la positivité que cet entrelaceur doit préserver.

\section{Principe d'Ambivalence : involution bidirectionnelle des rapports \(\pi/\varphi^2\) et \(\varphi/\pi^2\)}

On définit

\[
\mathcal Q(x,y)=\left(\frac{x}{y^2},\frac{y}{x^2}\right),
\qquad x,y>0.
\]

Pour

\[
P=xy,
\qquad
O=x/y,
\]

on vérifie

\[
P(\mathcal Q(x,y))=P^{-1},
\qquad
O(\mathcal Q(x,y))=O^3.
\]

Après deux étapes,

\begin{equation}
\boxed{P(\mathcal Q^2)=P,
\qquad O(\mathcal Q^2)=O^9.}
\label{eq:rm-qmap-nine}
\end{equation}

En coordonnées \(u=\log(xy)\), \(v=\log(x/y)\),

\[
(u,v)\longmapsto(-u,3v)\longmapsto(u,9v).
\]

Le produit revient et l'orientation multiplie sa mémoire par neuf. Pour
\((x,y)=(\pi,\varphi)\), les deux sorties sont

\[
\rho_A=\frac{\pi}{\varphi^2},
\qquad
\rho_E=\frac{\varphi}{\pi^2}.
\]

Elles satisfont

\[
\rho_E\rho_A^2=\varphi^{-3},
\qquad
M_9=(\rho_E\rho_A^2)^6=\varphi^{-18}.
\]

La relation relie la lecture aréale, l'exposant électronique, le retour nonadique
et la mémoire quadratique sans identifier leurs structures algébriques entre elles.

\begin{proofstatusbox}
Les algèbres quadratiques, le calcul dual, la réduction de Schur, la
courbure \eqref{eq:rm-trit-cartan-curvature} et l'identité
\eqref{eq:rm-qmap-nine} sont exactes. La limite de l'énergie torsionnelle et sa
réalisation comme action complète d'Einstein--Cartan exigent les
carrés de raffinement et le morphisme physique développés ultérieurement.
\end{proofstatusbox}

\begin{falsifierbox}
Réfutent la clôture interne : un terme autre que
\(-\tau\ell^{-2}e\wedge e\) dans la courbure ; un coefficient nilpotent qui ne
soit pas \(D_\omega e\) ; un résidu dans l'identité de Feshbach--Schur ; la perte
de positivité de \(D_9^*WD_9\) ; ou l'échec de l'une des deux équations
de \eqref{eq:rm-qmap-nine}.
\end{falsifierbox}
